O monitoramento da ETc é uma prática que possui diversas abordagens e métodos, dependendo do ambiente e do nível de precisão desejado. A ETc pode ser estimada por meio de métodos diretos, como a pesagem de lisímetros, ou indiretos, como a equação de \textcite{bernardo_irrigacao2008} apresentada na seção anterior. Entre essas abordagens, esta pesquisa propõe uma aplicação web para o monitoramento da ETc em tempo real, combinando métodos indiretos para esse fim.

A seguir serão apresentados outros trabalhos com abordagens semelhantes, e como a presente pesquisa supera as limitações encontradas nesses estudos.

\section{Integração de dados e padrões web para modelos de ETo}

\textcite{Jianting_webeva2009} abordam o desenvolvimento de sistemas de disseminação de dados em tempo real baseados na aquisição de dados. Eles fazem uso de dados geoespaciais para estimar a ETo diária, apontando a necessidade de métodos de acesso públicos a esses dados. Isso contrasta com os métodos \textit{ad-hoc} utilizados nos serviços existentes com o mesmo propósito.

Os autores desenvolveram um sistema protótipo que utiliza padrões da \textit{Open Geospatial Consortium} (OGC) e o protocolo \textit{Network Data Access Protocol} (OPeNDAP) para publicar dados de estações meteorológicas. A arquitetura desse sistema incorpora diversas fontes de dados distribuídas e heterogêneas, como registros diários e horários de estações geoestacionárias. Além disso, imagens de satélite e resultados de modelos meteorológicos são utilizados para uma análise espacial. 

Através de serviços web construídos sobre essas fontes, o sistema prova a eficácia de uma abordagem distribuída. Isso permite tanto a utilização quanto o desenvolvimento de novos serviços e aplicações sem interferir nos já existentes.

O trabalho de \textcite{Jianting_webeva2009} promove várias vantagens em relação a interoperabilidade e padronização de dados. No entanto, esse método não considera medições \textit{in situ}, o que em cenários menores, como o de uma propriedade rural, pode ser uma limitação em questão de precisão e custo. Além disso, o sistema proposto não oferece uma interface para a visualização e análise dos dados, o que pode dificultar a interpretação dos resultados por usuários finais.

\section{Ferramentas web para determinação de ETo}

O estudo de \textcite{Fangming_webeva2021} foca no desenvolvimento e implementação de APIs em nuvem, chamadas ETWatch, para a geração de dados regionais da ETo. A pesquisa utiliza imagens de satélite e técnicas de sensoriamento remoto, integrando modelos climáticos e hidrológicos para estimar a evapotranspiração. As APIs baseadas em computação em nuvem são desenvolvidas para proporcionar acesso fácil e em tempo real aos dados de ETo.

A metodologia combina dados de sensoriamento remoto, medições \textit{in situ} e modelagem matemática para calcular a evapotranspiração. As APIs facilitam a disseminação desses dados para diversos usuários e aplicações. Entre as vantagens do sistema ETWatch, destacam-se o acesso em tempo real aos dados de ETo, permitindo uma resposta rápida para a gestão de recursos hídricos. Além disso, o sistema fornece uma interface web para monitorar os dados, facilitando o acesso e a leitura das informações obtidas.

Apesar das vantagens, \textcite{Fangming_webeva2021} identificam alguns pontos que necessitam de maior preocupação. A precisão das estimativas pode ser limitada pela qualidade e disponibilidade dos dados de entrada, especialmente em regiões com poucas medições \textit{in situ}. A complexidade de integrar modelos climáticos e hidrológicos pode apresentar desafios computacionais e de implementação. Além disso, há a necessidade de ajustar os modelos para responder a variações e mudanças climáticas contínuas.

Uma abordagem semelhante foi realizada por \textcite{Jose_webeva2011} com o desenvolvimento de uma aplicação \textit{desktop} chamada \textit{CropWaterUse}. Essa aplicação foi desenvolvida para estimar os tempos de irrigação baseados na estimativa de ETc. Assim como \textcite{Fangming_webeva2021}, \textcite{Jose_webeva2011} utilizaram medições e modelos matemáticos para estimar a ETo. Contudo, eles realizaram o cálculo da ETc com base na ETo e em valores de $Kc$ específicos para cada cultura, sendo estes inseridos manualmente pelo usuário na aplicação.

Contudo, a aplicação \textit{CropWaterUse} não oferece uma interface de acesso distribuído para a visualização dos dados, seja por meio de APIs ou nuvens de dados. Isso limita sua utilização em cenários remotos ou até mesmo restritos, como propriedades rurais.

\section{Desenvolvimento e aplicação de ferramentas para ETo}